% GENERATED_BY: oc_core_monograph_translation_draft_pipeline_v1.py
% AI_ASSISTED_TRANSLATION_DRAFT: TRUE
% THEOREM_REVIEW_STATUS: PENDING
% THEOREM_REVIEW_MODE: FORMAL_AI_GATE__CODEX_CLI_CHATGPT
% THEOREM_REVIEW_REVIEWER_ID:
% THEOREM_REVIEW_AUTHORITY_CLASS:
% THEOREM_REVIEWED_AT:
% THEOREM_REVIEW_DOSSIER_REF:
% SOURCE_LANGUAGE: EN
% TARGET_LANGUAGE: RU
% SOURCE_EN_REF: content/m_spaces/m7.tex
% ================================================================
% ==== FILE: content/m_spaces/m7.tex
% ================================================================

% ==============================
%  Ontology of Continua — Core
%  Meta-Space: M7
%  FULL MODULE — FINAL VERSION
% ==============================

\subsubsection{\texorpdfstring{$M_7$}{M_7} Обзор}
\label{sec:m7-overview}

$M_7$ — это мета-пространство, которое делает допустимой \emph{социальную организацию}.
Хотя $M_6$ поддерживает когнитивные континуумы, основанные на аттракторах, понятиях и
протологических операциях, $M_7$ позволяет:

\begin{itemize}
    \item общий смысл и интерсубъективные понятия,
    \item социальные оси и механизмы координации,
    \item коммуникационные каналы и символический обмен,
    \item появление циклов и норм на уровне группы,
    \item потенциалы и потоки на уровне популяции,
    \item структурные условия существования социальных континуумов ($K_7$).
\end{itemize}

Следовательно, $M_7$ — минимальная среда для \emph{многоагентного когнитивного сопряжения}.

% ---------------------------------------------------------------
\subsubsection*{1. Допустимое пространство конфигураций \texorpdfstring{$\Omega(M_7)$}{\Omega(M_7)}}

\[
\Omega(M_7)=
\left\{
\big(
\mathcal{C}_i,\
A_{\mathrm{shared}},\
R_{\mathrm{comm}},\
\mathcal{N},\
C_{\mathrm{soc}},\
\Sigma_7
\big)
\ \big| \ \text{выполняются условия допустимости}
\right\}.
\]

Компоненты:

\begin{itemize}
    \item $\mathcal{C}_i$ — когнитивные континуумы отдельных индивидов, вложенные в $M_6$,
    \item $A_{\mathrm{shared}}$ — оси, поддерживающие общую репрезентацию и смысл,
    \item $R_{\mathrm{comm}}$ — допустимые коммуникационные отношения,
    \item $\mathcal{N}$ — сеть структуры взаимодействий,
    \item $C_{\mathrm{soc}}$ — циклы поддержания социального порядка (нормы, роли, взаимность),
    \item $\Sigma_7$ — ограничения когерентности, обеспечивающие устойчивость группы.
\end{itemize}

Допустимость требует:

\begin{enumerate}
    \item существования хотя бы одной общей семантической оси между агентами,
    \item возможности передачи различий вдоль $R_{\mathrm{comm}}$,
    \item устойчивости социальных циклов к флуктуациям,
    \item топологии взаимодействий $\mathcal{N}$, поддерживающей когерентность, а не фрагментацию.
\end{enumerate}

% ---------------------------------------------------------------
\subsubsection*{2. Допустимые оси \texorpdfstring{$A(M_7)$}{A(M_7)}}

Новые оси, вводимые в $M_7$:

\begin{itemize}
    \item $A_{\mathrm{shared}}$ — оси общего смысла / интерсубъективности,
    \item $A_{\mathrm{comm}}$ — коммуникационные оси (символы, сигналы, нормы),
    \item $A_{\mathrm{coord}}$ — оси координированных действий,
    \item $A_{\mathrm{role}}$ — оси, соответствующие социальным ролям,
    \item $A_{\mathrm{norm}}$ — оси, связанные с нормами и соблюдением правил,
    \item $A_{\mathrm{coop}}$ — спектр сотрудничества и дефекции,
    \item $A_{\mathrm{identity}}$ — оси, отражающие различия групповой идентичности.
\end{itemize}

Необходимое условие для $K_7$:

\[
\dim(A_{\mathrm{shared}})\ge 1,\
\quad
A_{\mathrm{comm}}\subseteq A(M_7).
\]

% ---------------------------------------------------------------
\subsubsection*{3. Потенциалы \texorpdfstring{$P(M_7)$}{P(M_7)}}

Социальные потенциалы обобщают $P(M_6)$ до многоагентной организации:

\[
P(M_7)=
(F_{\mathrm{comm}},\
F_{\mathrm{coord}},\
F_{\mathrm{status}},\
F_{\mathrm{norm}},\
F_{\mathrm{coh}},\
F_{\mathrm{coop}}).
\]

Интерпретация:

\begin{itemize}
    \item $F_{\mathrm{comm}}$ — потенциал коммуникации, обеспечивающий информационный поток,
    \item $F_{\mathrm{coord}}$ — потенциал, стабилизирующий согласованные действия,
    \item $F_{\mathrm{status}}$ — градиенты социального влияния и иерархии,
    \item $F_{\mathrm{norm}}$ — нормативные потенциалы (давление к соответствию),
    \item $F_{\mathrm{coh}}$ — потенциал групповой когерентности,
    \item $F_{\mathrm{coop}}$ — энергетический эквивалент стимулов кооперации.
\end{itemize}

Социальный континуум существует, когда:

\[
F_{\mathrm{coh}} > \Theta_{\mathrm{fragment}},
\quad F_{\mathrm{norm}} \text{ допускает устойчивые циклы}.
\]

% ---------------------------------------------------------------
\subsubsection*{4. Пороги \texorpdfstring{$\Theta(M_7)$}{\Theta(M_7)}}

Критические пороги включают:

\begin{itemize}
    \item $\Theta_{\mathrm{comm}}$ — минимальная пропускная способность коммуникации,
    \item $\Theta_{\mathrm{shared}}$ — порог формирования общего понятия,
    \item $\Theta_{\mathrm{coh}}$ — порог групповой когерентности,
    \item $\Theta_{\mathrm{coop}}$ — порог сотрудничества,
    \item $\Theta_{\mathrm{norm}}$ — минимальное давление для устойчивости норм,
    \item $\Theta_{\mathrm{fragment}}$ — порог фрагментации на $\partial\Omega(M_7)$.
\end{itemize}

Переход за $\Theta_{\mathrm{fragment}}$ приводит к коллапсу в независимые континуумы $K_6$.

% ---------------------------------------------------------------
\subsubsection*{5. Потоки \texorpdfstring{$J(M_7)$}{J(M_7)}}

Допустимые потоки:

\[
J(M_7)=
\left(
J_{\mathrm{comm}},\
J_{\mathrm{norm}},\
J_{\mathrm{coop}},\
J_{\mathrm{identity}},\
J_{\mathrm{coord}},\
\nabla_i T_7
\right).
\]

Интерпретация:

\begin{itemize}
    \item $J_{\mathrm{comm}}$ — коммуникационный поток по сети,
    \item $J_{\mathrm{norm}}$ — распространение норм,
    \item $J_{\mathrm{coop}}$ — динамика кооперации/дефекции,
    \item $J_{\mathrm{identity}}$ — потоки, формирующие групповую идентичность,
    \item $J_{\mathrm{coord}}$ — потоки, обеспечивающие синхронизированное действие,
    \item $\nabla_i T_7$ — градиент социальной напряженности.
\end{itemize}

Условие сохранения (ограниченные ресурсы внимания, доверия, вовлечения):

\[
\sum_i J_{\mathrm{comm},i} \le C_{\mathrm{channel}}.
\]

% ---------------------------------------------------------------
\subsubsection*{6. Циклы \texorpdfstring{$C(M_7)$}{C(M_7)}}

Фундаментальные социальные циклы:

\[
C_{\mathrm{soc}} =
\{
\text{сигнал} \to \text{ответ} \to \text{обратная связь} \to
\text{обновление норм} \to \text{новый сигнал}
\}.
\]

Другие циклы:

\begin{itemize}
    \item $C_{\mathrm{role}}$ — приобретение роли, реализация, подкрепление,
    \item $C_{\mathrm{coop}}$ — цикл сотрудничества (выгода $\to$ доверие $\to$ кооперация),
    \item $C_{\mathrm{conflict}}$ — конфликт $\to$ переговоры $\to$ разрешение,
    \item $C_{\mathrm{identity}}$ — формирование и стабилизация групповой идентичности,
    \item $C_{\mathrm{institution}}$ — возникновение институциональных правил.
\end{itemize}

Живой социальный континуум требует:

\[
\oint_{C_{\mathrm{soc}}} dA_{\mathrm{state}} \approx 0.
\]

% ---------------------------------------------------------------
\subsubsection*{7. Граница \texorpdfstring{$\partial\Omega(M_7)$}{\partial\Omega(M_7)}}

Конфигурация достигает $\partial\Omega(M_7)$, когда:

\begin{itemize}
    \item общие оси схлопываются (утрата общего смысла),
    \item коммуникационные каналы выходят из строя или насыщаются,
    \item нормативные циклы разрушаются (нет устойчивых ожиданий),
    \item кооперация становится неустойчивой,
    \item оси идентичности распадаются на несогласованные подпространства,
    \item социальное напряжение $T_7$ превышает порог допустимости.
\end{itemize}

Пересечение границы приводит к фрагментации на несколько континуумов $K_6$.

% ---------------------------------------------------------------
\subsubsection*{8. Связь с \texorpdfstring{$M_6$}{M_6} и \texorpdfstring{$M_8$}{M_8}}

\textbf{От $M_6$ к $M_7$:}

Необходимые условия:

\[
A_{\mathrm{concept}}^{(i)} \cap A_{\mathrm{concept}}^{(j)} \neq \emptyset,
\quad
R_{\mathrm{comm}} \neq \emptyset.
\]

Общая концептуальная структура требуется для интерсубъективности.

\textbf{К $M_8$ (цивилизационно-технологическое пространство):}

$M_8$ вводит:

\begin{itemize}
    \item символические системы,
    \item письмо, институты, инфраструктуру,
    \item технологические и экономические оси.
\end{itemize}

Поэтому $M_7$ обеспечивает основание \emph{протоинституциональной когерентности}, необходимой для $K_8$.
